\centerline{\Huge \bf  Олимпиада II}
\centerline{\Huge \bf 7 класс}

\begin{flushright}
    \scriptsize{Заключительный отбор в Сириус 2022г.\\ Кто загуглит - тому бан!}
\end{flushright}

\problem{1}
Четыре вершины куба, образующие его нижнюю грань, окрашены в красный, зеленый, синий и оранжевый цвета. Сколько существует способов раскрасить четыре верхних вершины куба в эти цвета так, чтобы каждая грань имела четыре вершины разного цвета? Ответ объясните.

\problem{2}
На шести карточках написаны цифры $1, 2, 3, 4, 5, 6$. Составьте из них шестизначное число $\overline{abcdef}$ такое, что $\overline{ab}$ делится на $2$, $\overline{abc}$ делится на $3$, $\overline{abcd}$ делится на $4$, $\overline{abcde}$ делится на $5$, $\overline{abcdef}$ делится на $6$. Сколько таких чисел можно получить?


\problem{3}
Во время подготовки к очному туру Миша усердно решал задачи из доступных ему курсов Сириуса. В некоторых задачах он мог себя проверить по ответу. В остальных задачах ответа не было. Число задач с ответом составляет $9/11$ от числа остальных. Миша решил $2/7$ задач с ответом и $3/5$ остальных. В итоге ему осталось решить больше $500$ задач с ответом и меньше $500$ остальных. Сколько задач было в доступных Мише курсах Сириуса?

\problem{4}
Найдите все пары простых чисел $(p, q)$, для которых $15(p - 1)(q - 1)$ делится на $pq$.

\problem{5}
В треугольнике $ABC$ угол $A$ в $3$ раза больше угла $B$, который в $2$ раза больше угла $C$. Сторона $BC$ на $5$ длиннее стороны $AC$. Найдите длину биссектрисы треугольника, проведенной из вершины $A$.

\problem{6}
Назовём число $2022$ замечательным. Известно, что для любого целого m в множестве $\{m, 2m + 1, 3m\}$, все три числа замечательны, если замечательно хотя бы одно из них. Можно ли утверждать, что число $20222022$ замечательно?